\scanpage{146}
évident, d'ailleurs pleinement fidèle
\[
 \underline h^{\deleted{\unsure{\circ|0|\#}}}:\mathscr V\longrightarrow\mathcal M^{\deleted{\circ}}
\]
en posant $\underline h^{\deleted{\circ}}(X)=(X,\mathrm{id}_X)$. Pour toute catégorie abélienne\inserted{ $\mathcal C$},
% slanted
\margin{il suffit que $\mathcal C$ soit une catégorie additive où les projecteurs sont réalisables}
le foncteur $\varphi:T\mapsto T\circ h$
\[
 \underline{\operatorname{Fonc\,add}}(\mathcal M,\mathcal C)\xrightarrow{\varphi}\underline{\operatorname{Fonc\,add}}(\mathscr V,\mathcal C)
\]
est une \deleted{isomorphisme} équivalence. Pour le voir, on va construire un foncteur quasi-inverse \deleted{1°) pleinement fidèle. Un}
\[
 \psi:\underline{\operatorname{Fonc\,add}}(\mathscr V,\mathcal C)\longrightarrow\underline{\operatorname{Fonc\,add}}(\mathcal M,\mathcal C)
\]
en associant à \unsure{tout|tt} foncteur\inserted{ additif} $H:\mathscr V\unsure{\mathord{-}|\to}\mathcal C$ son extension canonique
\[
 \psi(H)(X,\pi)=\operatorname{Im}H(\pi)
\]
avec la loi fonctorielle évidente. On voit aussitôt que l'on a des isom can
\[
\begin{aligned}
 \varphi\deleted{\unsure{(}}\psi\deleted{\unsure{)}}(H)&\simeq H\\
 \psi\varphi(T)&\simeq T
\end{aligned}
\]

Pour l'instant, cette construction et propriété universelle de $\mathcal M$ n'utilise pas les hyp particulières faites sur $\mathscr V$ ; elle répond au Pb universel de rendre \uline{tous les projecteurs réalisables}.\inserted{ (ds catégorie additive donnée)} Mais les hypothèses faites sur $\mathscr V$ vont impliquer que $\mathcal M$ est en fait une catégorie \inserted{\uline{abélienne}} \uline{semi-simple} (i.e. tout objet de
% archive sheet number: 145
